\documentclass{article}
\usepackage{amsmath,amssymb}
\usepackage{hyperref}
% Shared commands for notes in docs/paper-gaps/.

\DeclareUnicodeCharacter{2080}{\ifmmode{}_{0}\else\textsubscript{0}\fi}
\DeclareUnicodeCharacter{2081}{\ifmmode{}_{1}\else\textsubscript{1}\fi}
\DeclareUnicodeCharacter{2082}{\ifmmode{}_{2}\else\textsubscript{2}\fi}
\DeclareUnicodeCharacter{2099}{\ifmmode{}_{n}\else\textsubscript{n}\fi}

\newcommand{\Lean}{\textsc{Lean}}
\ExplSyntaxOn
\NewDocumentCommand{\leanid}{m}{
  \ifmmode
    \textsf{\detokenize{#1}}
  \else
    \group_begin:
    \urlstyle{sf}
    \tl_set:Nn \l_tmpa_tl {#1}
    \tl_replace_all:Nnn \l_tmpa_tl {\_} {_}
    \exp_args:NV \path \l_tmpa_tl
    \group_end:
  \fi
}
\ExplSyntaxOff
\providecommand{\tr}{\operatorname{tr}}
\newcommand{\tnleanIssueBase}{https://github.com/LionSR/TNLean/issues/}
\newcommand{\tnleanPrBase}{https://github.com/LionSR/TNLean/pull/}
\newcommand{\tnleanDocBase}{https://sirui-lu.com/TNLean/docs/}
\newcommand{\ghissue}[1]{\href{\tnleanIssueBase#1}{GitHub issue \##1}}
\newcommand{\ghpr}[1]{\href{\tnleanPrBase#1}{PR \##1}}
\newcommand{\doclink}[1]{\href{\tnleanDocBase#1.html}{\path{#1}}}

% Dirac notation and the identity operator, as in blueprint/src/macros/common.tex.
% \providecommand keeps note-local definitions authoritative.
\usepackage{dsfont}
\providecommand{\ket}[1]{|#1\rangle}
\providecommand{\bra}[1]{\langle #1|}
\providecommand{\braket}[2]{\langle #1 | #2 \rangle}
\providecommand{\ketbra}[2]{|#1\rangle\!\langle #2|}
\providecommand{\Id}{\mathds{1}}
\providecommand{\one}{\mathds{1}}
\providecommand{\C}{\mathbb{C}}
\providecommand{\MN}[1]{M_{#1}(\C)}
\providecommand{\MD}{M_D(\C)}

% Verdict marker: \gapnote{<kind>}{<status>}. Kind is the policy
% classification (clarification, local-correction, scope-restriction,
% unfaithful, false-source, open-gap); status is open, wip (elimination
% underway), resolved, or historical. Machine-read by the site index;
% prints a verdict line.
\newcommand{\gapnote}[2]{\par\noindent\textbf{Verdict:} #1 (#2).\par}

\title{Critical Parameters and Boundary Localization in the One-Species PVBS Chain}
\author{}
\date{2026-10-04}
\begin{document}
\maketitle
\gapnote{false-source}{resolved}

\section{The assertion}
The review arXiv:2011.12127, Section IV.C.4 (HTML Section IV.3.4), gives the
matrices $A^0=\operatorname{diag}(1+\lambda,1)$ and
$A^1=|0\rangle\langle1|$, with $\lambda\in[-1,1]$.
It describes a two-dimensional open-chain ground space, a boundary-localized
particle, and a unique product vacuum on the periodic chain. The printed right-edge
condition is $\lambda>1$.

\section{The point at issue}
Write $q=1+\lambda$. A particle at the zero-indexed site $k$ has amplitude
$q^k$. Thus left localization occurs for $0\le q<1$ and right localization
for $q>1$: in the review's parameter these are $\lambda<0$ and $\lambda>0$.
The printed right-edge condition is outside its stated parameter interval.
At $q=1$, the particle is the delocalized W vector. On a two-site ring,
$|10\rangle+|01\rangle$ satisfies both local constraints and is independent
of the vacuum. Therefore unconditional periodic vacuum uniqueness is false.

\section{The corrected statement}
For any complex $q$, the open parent Hamiltonian with two-site interactions
has ground space
\begin{align}
    \operatorname{span}\left\{|0\cdots0\rangle,
    \sum_{k=0}^{N-1}q^k|0\cdots1_k\cdots0\rangle\right\},\qquad N\ge2.
    \label{eq:pvbs_gap_open_space}
\end{align}
Its dimension is two, including $q=0$ and $q=1$.
On a periodic chain, the vacuum is the unique ground state whenever $q^N\ne1$.
In particular, this holds at every length $N\ge2$ if $|q|\ne1$.
The latter condition is sufficient, not necessary at any fixed length.
To prove uniqueness, write any periodic ground vector as vacuum plus $b$ times
the open-chain particle. A one-site translation remains a ground vector.
Comparing the first two particle amplitudes forces $b=bq^N$, hence $b=0$.
The vacuum satisfies every local constraint.

For $|q|<1$, the left-edge amplitudes decay as $|q|^k$. For $|q|>1$, the
amplitude $k$ sites from the right edge, divided by the rightmost amplitude,
is $q^{-k}$. These describe localization without asserting a uniform spectral gap.

\section{Source conventions and conclusion}
Bachmann and Nachtergaele, arXiv:1112.4097, Section II, equations (1)--(8),
use positive hopping parameters and phases, and explicitly exclude unit hopping
in their boundary-localization discussion. Their matrix-product ordering differs
from the left-to-right ordering above; the displayed one-particle amplitudes
fix the convention. More explicitly, $q=\mu e^{i\theta}$ and the forbidden
local ket is $|01\rangle-\overline q|10\rangle$, whose orthogonal
constraint is $a_{01}=q a_{10}$. The primary source assumes $\mu>0$;
$q=0$ is an algebraic endpoint extension. The formalization permits complex $q$, while the review's
real family is recovered by $q=1+\lambda$.
The corrected finite-chain ground spaces and the critical two-site counterexample
resolve the discrepancy. No spectral-gap theorem is claimed here.\footnote{Formal declarations:
    \leanid{MPSTensor.ker_openParentHamiltonianES_pvbs},
    \leanid{MPSTensor.ker_parentHamiltonian_pvbs},
    \leanid{MPSTensor.pvbs_critical_periodic_counterexample}, and
    \leanid{MPSTensor.pvbsEdge_right_decay}.}
\end{document}
